% EVERY CANDIDATE BREAK IS WEIGHED AGAINST ITS OWN LINE'S WIDTH.
%
% A break that would start line five is measured against the width line
% five will have, not against whatever line the first active node
% happens to sit on. hstex used one width for every candidate at a
% breakpoint, which is right only while all the lines are the same
% width -- and wrong the moment \hangindent or \parshape makes them
% differ.
%
% With \hangindent=30pt \hangafter=1 over a paragraph of short words,
% \tracingparagraphs shows it at once:
%
%   ref:    @ via @@1 b=6396 p=0 d=41036836
%   hstex:  @ via @@1 b=10000 p=0 d=100000000
%
% -- the candidate from @@1 would begin a hung line thirty points
% narrower, and measuring it against the full width made every such
% break look infinitely bad. The paragraph then broke into fewer lines,
% and the feasible breaks the reference keeps -- one active node per
% line number where the widths differ -- never appeared at all.
%
% The last paragraph here has no hanging and is the control: it agreed
% before the fix and after it.
%
% STILL DIFFERING, in the \parshape case only: the reference keeps ONE
% ACTIVE NODE PER LINE NUMBER where the widths differ, and hstex keeps
% one per fitness class however the line numbers fall. So at one
% breakpoint the reference records
%
%   @@2: line 1.0 t=100010000 -> @@0
%   @@3: line 2.0 t=200010000 -> @@1
%
% and hstex records only the first. The line numbers stop mattering once
% they are past the last line \parshape or \hangindent names -- ten pairs
% of \parshape make lines eleven onward all one class -- so the bucket is
% (fitness, line) up to that point and (fitness) after it.
%
% ATTEMPTED AND BACKED OUT. Bucketing the candidates by class and making
% one set of active nodes per class is not enough on its own, because it
% rests on THE ACTIVE LIST BEING IN LINE ORDER and hstex's is not. The
% reference walks the list in order, and when it reaches a node whose
% line number is past the class it was gathering it makes that class's
% nodes and puts them IN AT THAT POINT -- which is both what keeps the
% list sorted and what puts the `@@n:' lines between the `@ via' lines.
% hstex appends new records at the end, so a class can be met twice and
% the nodes come out in the wrong order:
%
%   @@3: line 2.0 t=200010000 -> @@2      hstex, decreasing
%   @@4: line 1.0 t=100010000 -> @@0
%
% Doing this properly means inserting an active node at the walk's
% position rather than at the end.
%
% WRITTEN A SECOND TIME AND BACKED OUT AGAIN. The insertion is easy --
% rebuild the list into a second buffer and swap -- but the CLASS RULE is
% not what it looks like, and getting it wrong breaks cases that were
% right. Eight measurements, all `does one breakpoint make one record or
% two when the candidates are line 1 and line 2':
%
%   \hangafter=0                    ONE
%   \hangafter=1                    ONE
%   \hangafter=2                    TWO
%   \hangafter=3                    TWO
%   \hangafter=-3                   TWO
%   \parshape=1                     ONE
%   \parshape=3                     TWO
%   \looseness=1, nothing else      TWO
%
% Write E for the last line whose width is its own -- \parshape's pair
% count, or |\hangafter|, or nought -- except that any \looseness at all
% makes E unbounded. Then the only rule that fits all eight is
%
%   class(line) = min(line, E)
%
% and NOT `line up to E, one shared class after it', which is what the
% widths would suggest and what was tried. The difference shows exactly
% at E <= 1, where min() puts lines 1 and 2 together and the widths do
% not. Check it once more before building on it: find a breakpoint whose
% candidates are lines 2 and 3 under \hangafter=2, where min() says one
% class and the widths say two.
%
% Also worth knowing: when two classes DO get records, the reference
% writes them interleaved -- `@ via @@0', the class's `@@n:', `@ via
% @@1', its `@@n:' -- and when they do not, both `@ via' lines come
% first and one `@@n:' follows. That is a quick way to read which
% happened.
\tracingonline=1 \hsize=100pt \parindent=0pt \tracingparagraphs=1
\pretolerance=-1 \tolerance=10000 \hbadness=10000
\message{[A hangindent, so lines differ in width]}
\hangindent=30pt \hangafter=1
\noindent aa bb cc dd ee ff gg hh ii jj kk ll mm nn oo pp\par
\message{[B parshape]}
\parshape=3 0pt 40pt 10pt 60pt 0pt 100pt
\noindent aa bb cc dd ee ff gg hh ii jj kk ll mm nn oo pp\par
\message{[C no hanging, the control]}
\hangindent=0pt
\noindent aa bb cc dd ee ff gg hh ii jj kk ll mm nn oo pp\par
\message{[done]}
\end
